% GENERATED FILE. Edit canonical lessons, then run npm run generate:books.
% canonical-source-hash: fnv1a64:1abe7ecfa04d2a86
% canonical-lessons: FR-C290-crime, FR-C290-criminel, FR-C290-meurtre, FR-C290-delinquant, FR-C290-delinquance

\chapter{Le crime}
\label{ch:fr-le-crime}
\hlchaptermodality{290}

\begin{chapteropening}
\textbf{By the end of this chapter:} \emph{I can talk about serious crime: a crime, a criminal, a murder, a young offender and crime in general.}

Complete the last lesson of chapter 290: I can talk about serious crime: a crime, a criminal, a murder, a young offender and crime in general.
\end{chapteropening}

\section[le crime]{le crime — the crime; the murder}

\label{lesson:FR-C290-crime}



\begin{quote}
Before the new one: say the French for to threaten, then the French for the threat.
\end{quote}

\subsection*{You'll want to know: le crime}
\textbf{le crime} — \textquotedblleft{}the crime; the murder\textquotedblright{}.

\begin{quote}
Ce crime a choqué toute la ville.
\end{quote}

\begin{quote}
\emph{This crime shocked the whole town.}
\end{quote}

In French, \emph{un crime} is a serious offence, often a murder.

\begin{cousinweb}
\emph{crime} is inherited from Latin \emph{crīmen} \textquotedblleft{}a charge, a crime\textquotedblright{} (EtymDB).
\end{cousinweb}

\begin{grammarlens}[title={What you've built}]
One more word for telling, explaining and describing at B1.
\end{grammarlens}

\subsection*{Your turn}
\begin{itemize}
\raggedright
  \item \emph{Say it:} \emph{le crime}
  \item \emph{Say it:} \emph{le crime}, once more
  \item \emph{Say it:} a sentence of your own with \emph{le crime}
  \item \emph{Recall:} say the French for to threaten, then the French for the threat, then say \emph{le crime} again
\end{itemize}

\begin{tcolorbox}[breakable,colback=teal!4,colframe=teal!35!black,title={Before you move on}]
What does le crime mean? (\textquotedblleft{}The crime; the murder\textquotedblright{}.) Say it once more.
\end{tcolorbox}
\section[le criminel]{le criminel — the criminal}

\label{lesson:FR-C290-criminel}



\begin{quote}
Before the new one: say the French for the threat, then the French for the crime.
\end{quote}

\subsection*{You'll want to know: le criminel}
\textbf{le criminel} — \textquotedblleft{}the criminal\textquotedblright{}.

\begin{quote}
Le criminel a été arrêté deux jours plus tard.
\end{quote}

\begin{quote}
\emph{The criminal was arrested two days later.}
\end{quote}

Also an adjective: \emph{un acte criminel}.

\begin{cousinweb}
\emph{criminel} is a learned borrowing from Latin \emph{crīminālis} (EtymDB), built on \emph{crīmen} \textquotedblleft{}a charge, a crime\textquotedblright{}.
\end{cousinweb}

\begin{grammarlens}[title={What you've built}]
One more word for telling, explaining and describing at B1.
\end{grammarlens}

\subsection*{Your turn}
\begin{itemize}
\raggedright
  \item \emph{Say it:} \emph{le criminel}
  \item \emph{Say it:} \emph{le criminel}, once more
  \item \emph{Say it:} a sentence of your own with \emph{le criminel}
  \item \emph{Recall:} say the French for the threat, then the French for the crime, then say \emph{le criminel} again
\end{itemize}

\begin{tcolorbox}[breakable,colback=teal!4,colframe=teal!35!black,title={Before you move on}]
What does le criminel mean? (\textquotedblleft{}The criminal\textquotedblright{}.) Say it once more.
\end{tcolorbox}
\section[le meurtre]{le meurtre — the murder}

\label{lesson:FR-C290-meurtre}



\begin{quote}
Before the new one: say the French for the crime, then the French for the criminal.
\end{quote}

\subsection*{You'll want to know: le meurtre}
\textbf{le meurtre} — \textquotedblleft{}the murder\textquotedblright{}.

\begin{quote}
Le roman raconte l'histoire d'un meurtre.
\end{quote}

\begin{quote}
\emph{The novel tells the story of a murder.}
\end{quote}

The killer is \emph{le meurtrier}, \emph{la meurtrière}.

\begin{cousinweb}
\emph{meurtre} is built on Old French \emph{murtrir} \textquotedblleft{}to murder\textquotedblright{}, which is usually traced to a Frankish verb related to English \emph{murder} (EtymDB lists no parent).
\end{cousinweb}

\begin{grammarlens}[title={What you've built}]
One more word for telling, explaining and describing at B1.
\end{grammarlens}

\subsection*{Your turn}
\begin{itemize}
\raggedright
  \item \emph{Say it:} \emph{le meurtre}
  \item \emph{Say it:} \emph{le meurtre}, once more
  \item \emph{Say it:} a sentence of your own with \emph{le meurtre}
  \item \emph{Recall:} say the French for the crime, then the French for the criminal, then say \emph{le meurtre} again
\end{itemize}

\begin{tcolorbox}[breakable,colback=teal!4,colframe=teal!35!black,title={Before you move on}]
What does le meurtre mean? (\textquotedblleft{}The murder\textquotedblright{}.) Say it once more.
\end{tcolorbox}
\section[le délinquant]{le délinquant — the offender, the delinquent}

\label{lesson:FR-C290-delinquant}



\begin{quote}
Before the new one: say the French for the criminal, then the French for the murder.
\end{quote}

\subsection*{You'll want to know: le délinquant}
\textbf{le délinquant} — \textquotedblleft{}the offender, the delinquent\textquotedblright{}.

\begin{quote}
Le juge a parlé avec le jeune délinquant.
\end{quote}

\begin{quote}
\emph{The judge spoke with the young offender.}
\end{quote}

\emph{Un délinquant mineur} — a young offender.

\begin{cousinweb}
\emph{délinquant} is a learned word from Latin \emph{dēlinquēns}, the present participle of \emph{dēlinquere} \textquotedblleft{}to fail, to do wrong\textquotedblright{} (Lewis \& Short).
\end{cousinweb}

\begin{grammarlens}[title={What you've built}]
One more word for telling, explaining and describing at B1.
\end{grammarlens}

\subsection*{Your turn}
\begin{itemize}
\raggedright
  \item \emph{Say it:} \emph{le délinquant}
  \item \emph{Say it:} \emph{le délinquant}, once more
  \item \emph{Say it:} a sentence of your own with \emph{le délinquant}
  \item \emph{Recall:} say the French for the criminal, then the French for the murder, then say \emph{le délinquant} again
\end{itemize}

\begin{tcolorbox}[breakable,colback=teal!4,colframe=teal!35!black,title={Before you move on}]
What does le délinquant mean? (\textquotedblleft{}The offender, the delinquent\textquotedblright{}.) Say it once more.
\end{tcolorbox}
\section[la délinquance]{la délinquance — crime, delinquency}

\label{lesson:FR-C290-delinquance}



\begin{quote}
Before the new one: say the French for the murder, then the French for the offender.
\end{quote}

\subsection*{You'll want to know: la délinquance}
\textbf{la délinquance} — \textquotedblleft{}crime, delinquency\textquotedblright{}.

\begin{quote}
La délinquance a baissé dans le quartier.
\end{quote}

\begin{quote}
\emph{Crime has gone down in the neighbourhood.}
\end{quote}

Crime in general, as a problem of society.

\begin{cousinweb}
\emph{délinquance} is a learned borrowing from Late Latin \emph{dēlinquentia} (EtymDB), from \emph{dēlinquere} \textquotedblleft{}to do wrong\textquotedblright{}.
\end{cousinweb}

\begin{grammarlens}[title={What you've built}]
One more word for telling, explaining and describing at B1.
\end{grammarlens}

\subsection*{Your turn}
\begin{itemize}
\raggedright
  \item \emph{Say it:} \emph{la délinquance}
  \item \emph{Say it:} \emph{la délinquance}, once more
  \item \emph{Say it:} a sentence of your own with \emph{la délinquance}
  \item \emph{Recall:} say the French for the murder, then the French for the offender, then say \emph{la délinquance} again
\end{itemize}

\begin{tcolorbox}[breakable,colback=teal!4,colframe=teal!35!black,title={Before you move on}]
What does la délinquance mean? (\textquotedblleft{}Crime, delinquency\textquotedblright{}.) Say it once more.
\end{tcolorbox}
